\chapter{Server Setup \& Infrastruktur}

\section{Ziel und Aufbau}

Der Server bildet das zentrale Rückgrat der Smart-Home-Lernumgebung. Er stellt die Grundlage für Netzwerkzugang, Fernverwaltung, Containerisierung und zentrale Automationslogik bereit. In diesem Kapitel installierst und konfigurierst du die Infrastruktur schrittweise.

\Sidebar{Working Principle}
Ein Server muss in der Regel drei Dinge zuverlässig leisten: Zugriff auf das Internet, lokale Verteilung im Projekt-Netz und Remote-Verwaltung. Das Design legt dabei bewusst die Verantwortung auf den Xubuntu-Host, während der A1-Router nur noch als Netzwerkausgang und Access Point dient.

\section{Xubuntu Installation}

\begin{enumerate}
    \item \textbf{Stelle sicher, dass das Servergerät keine wichtigen Daten enthält, da die Installation das System vollständig überschreibt.}
    \item Lade das aktuelle \textbf{Xubuntu-24.04-Image} von der offiziellen Seite (\url{https://xubuntu.org/release/24.04/}) herunter und brenne es mit der Software \textbf{balenaEtcher} (\url{https://etcher.balena.io/}) auf einen USB-Stick. \KeyImage[]{Screenshot}{balenaEtcher: USB-Image auf Stick brennen}{../../media/balena-etcher.png}
    \item Schließe den USB-Stick an das Servergerät an und starte das System über das BIOS- bzw. Boot-Menü. Manche Hersteller verlangen ein ändern der Boot-Reihenfolge, sodass der USB-Stick als erstes Gerät geladen wird.
    \begin{InfoBox}{Wie gelange ich ins BIOS bzw. Boot-Menü?} 
        \scriptsize Je nach Hersteller und Modell des Geräts kann der Zugang unterschiedlich sein. Häufig sind die Tasten \texttt{F2}, \texttt{F10}, \texttt{F12} oder \texttt{ESC} relevant. Beim Neustart des Geräts achte auf die entsprechende Anzeige, die den Zugang zum BIOS oder Boot-Menü beschreibt und drücke die passende Taste rechtzeitig. \textbf{Alle Herstellerhinweise findet man frei zugänglich im Internet.}
    \end{InfoBox}\Sidebar{BIOS Menü}
    \item Wähle die Option \textit{Try or install Xubuntu}.
    \item Führe den Installationsprozess durch. Erstelle dabei einen passenden Benutzer mit einem sicheren Passwort und benenne das Gerät. Für den durchgeführten Testaufbau wählen wir als Benutzer das Kürzel der Schule \texttt{mmse} und als Gerätenamen \texttt{mmse-smarthome}.
\end{enumerate}

\begin{InfoBox}{Hinweis zur Xubuntu-Installation}
    \scriptsize Die Installation von Xubuntu ist ein einmaliger Schritt. Danach kann das System über die grafische Oberfläche oder per SSH verwaltet werden. Eine detaillierte Schritt-für-Schritt-Anleitung für die Installation von Xubuntu ist im Internet verfügbar, z. B. auf \url{https://ubuntu.com/desktop/docs/en/latest/tutorial/install-ubuntu-desktop/}.
\end{InfoBox}

\begin{InfoBox}{Hinweis zu alten Geräten}
    \scriptsize Sollte ein sehr altes Gerät verwendet werden, kann es vorkommen, dass bestimmte Komponenten wie WLAN-Chipsätze nicht mehr unterstützt werden. In diesem sehr seltenen Fall kann es notwendig sein, auf ein aktuelleres Gerät auszuweichen oder Ausweichlösungen zu nutzen.
    
    \textit{Der verwendete Toshiba Satellite Pro Laptop, wird mit seinem Broadcom-WLAN-Chipsatz nicht unterstützt. Hier wurde ein Edimax EW-7811UN USB-WLAN-Stick als Ersatz verwendet.}

    \textbf{In diesem Schritt solltest du jedenfalls eine temporäre Internetverbindung über LAN oder WLAN sicherstellen, um die Installation und spätere Updates durchführen zu können.} Diese Verbindung kann später wieder entfernt werden, sobald das lokale Projekt-Netzwerk korrekt eingerichtet ist.
\end{InfoBox}

\KeyImage[]{Screenshot}{Erfolgreicher Xubuntu Startup}{../../media/Xubuntu-Desktop.png}

\section{SSH einrichten}

SSH ermöglicht den sicheren Zugriff auf den Server über die Kommandozeile. Damit kannst du das System auch von einem anderen Rechner aus verwalten, ohne direkt am Gerät zu sitzen.

\KeyImage[width=0.6\textwidth]{Screenshot}{Terminal öffnen}{../../media/terminal-from-start-menu.png}

\CodeFile[SSH-Installation und Aktivierung auf dem Xubuntu-Server]{bash}{code/ssh-setup.sh}\label{lst:ssh-setup}

Die im Listing~\ref{lst:ssh-setup} (\nameref{lst:ssh-setup}) zu sehenden Befehle bewirken beim Eingeben in ein geöffnetes Terminal Folgendes:

\begin{enumerate}
    \item Installiert den aktuellsten OpenSSH-Server auf dem Xubuntu-Host.
    \item Aktiviert den Dienst automatisch beim Systemstart.
    \item Öffnet den SSH-Port in der Firewall und prüft den Verbindungsstatus.
    \item Gibt die lokale IP-Adresse des Servers aus, die für die Verbindung von einem anderen Rechner benötigt wird.
\end{enumerate}

\begin{InfoBox}{SSH-Verbindung testen}
    \scriptsize Teste die Verbindung von einem anderen Rechner im gleichen Netzwerk mit dem Befehl \texttt{ssh <Benutzer>@<IP-Adresse>}. Dabei ersetzt du \texttt{<Benutzer>} durch den zuvor erstellten Benutzernamen und \texttt{<IP-Adresse>} durch die vom Server ausgegebene Adresse. Bestätige die Sicherheitsabfrage mit \texttt{yes} und gib das Passwort ein, um die Verbindung herzustellen. Danach kannst du den Server über die Kommandozeile steuern.
\end{InfoBox}

\subsection{Passwortlose SSH Verbindung}

Sollte die Verbindung regelmäßig genutzt werden, kann es sinnvoll sein, den Rechner, mit dem man sich auf den Server verbindet, als vertrauenswürdiges Gerät zu hinterlegen. Dadurch \textbf{entfällt die Eingabe des Benutzerpassworts} bei jeder SSH-Verbindung. 
Dafür erstellen wir einmalig einen sog. \textit{SSH-Key} auf unserem Windows-Rechner und übertragen den öffentlichen Teil des Schlüssels auf den Server. Danach authentifiziert sich der Rechner ab der nächsten SSH-Verbindung automatisch durch seinen privaten Schlüssel, ohne dass wir das eigentliche Passwort eingeben müssen.

\KeyImage[]{Screenshot}{PowerShell in Windows öffnen}{../../media/open-powershell-windows.png}

\CodeFile[PowerShell-Befehle für SSH-Keys in Windows]{bash}{code/ssh-key-setup.ps1}

\begin{InfoBox}{Hinweis zu Benutzername und Hostname}
    \scriptsize Der Benutzername und Hostname des Servers müssen in den Befehlen angepasst werden.
\end{InfoBox}

\begin{InfoBox}{SSH-Keys für MacOS und Linux}
    \scriptsize Auf MacOS und Linux ist die Vorgehensweise ähnlich. Auch dort können SSH-Keys erstellt und auf den Server übertragen werden. Die genauen Befehle findest du unter \url{https://man.openbsd.org/ssh-keygen.1}.
    
\end{InfoBox}

\section{Remote Desktop konfigurieren}

Mit Remote Desktop kannst du zusätzlich zur Fernsteuerung über die Kommandozeile den grafischen Desktop des Xubuntu-Servers über das Netzwerk ansteuern. Das ist besonders hilfreich, wenn du GUI-basierte Konfigurationen oder Wartungsarbeiten durchführst, bei denen das grafische Einstellungsmenü hilfreich sein könnte.

\begin{enumerate}
    \item Terminal öffnen
    \item Aktualisiere das System und installiere anschließend den RDP-Server \texttt{xrdp}.
    \item Konfiguriere \texttt{xrdp} so, dass es die XFCE-Oberfläche verwendet. \Sidebar{Xubuntu basiert auf XFCE}
    \item Vergib die nötigen Zertifikatsrechte und starte den Dienst neu.
    \item Öffne den RDP-Port in der Firewall und prüfe mit \texttt{hostname -I} die Serveradresse.
\end{enumerate}

\CodeFile[Xubuntu-Setup für Remote Desktop]{bash}{code/rdp-setup.sh}

\Sidebar{Vor der Remote Desktop Verbindung}
Wenn der Laptop lokal weiter aktiv ist, kann der gleiche Benutzer nicht gleichzeitig per Remote Desktop angemeldet sein. \textbf{Melde dich deshalb immer erst lokal ab, bevor du eine neue Remote-Verbindung startest.} Das verhindert schwarze Bildschirme und Abbrüche.

Verwende für die Remote-Verbindung den Windows-Dienst \emph{Remotedesktopverbindung} und trage dort den Hostnamen oder die IP-Adresse des Servers ein. Danach kannst du dich mit dem zuvor erstellten Benutzer anmelden.

\KeyImage[width=0.6\textwidth]{Screenshot}{Windows-Remote-Desktop}{../../media/Remote-Desktop-Verbindung.png}

\subsection{Energieoptionen}

Falls es sich beim Servergerät um einen Laptop handelt, müssen die systemweiten Energieeinstellungen angepasst werden. Dadurch wird verhindert, dass das Gerät nach inaktiver Zeit oder bei geschlossenem Deckel in den Energiesparmodus wechselt. Das ist besonders wichtig, wenn \textbf{kein Benutzer lokal angemeldet ist}, da sonst die USB-Stromversorgung (z.\,B. für den WLAN-Stick) gekappt wird.

\CodeFile[Energieoptionen für Laptop-Server]{bash}{code/energieoptionen.sh}\label{lst:energieoptionen}

Die Anpassung erfolgt wie in Listing~\ref{lst:energieoptionen} (\nameref{lst:energieoptionen}) in zwei Schritten:

\begin{enumerate}
    \item \textbf{Aktion beim Schließen des Deckels (Lid Switch) deaktivieren:}\\
    Öffne die Systemd-Login-Konfiguration im Terminal mit \texttt{sudo nano /etc/systemd/logind.conf}. Suche dort nach den folgenden Zeilen, \textbf{entferne das \texttt{\#}-Zeichen am Zeilenanfang} und setze den Wert jeweils auf \texttt{ignore}:
    \begin{itemize}
        \item \texttt{HandleLidSwitch=ignore}
        \item \texttt{HandleLidSwitchExternalPower=ignore}
        \item \texttt{HandleLidSwitchDocked=ignore}
    \end{itemize}
    Speichere die Datei (\texttt{Strg+O}, \texttt{Enter}, \texttt{Strg+X}) und starte den Dienst mit \texttt{sudo systemctl restart systemd-logind} neu (\textit{kann die Session ggf. kurz unterbrechen}).

    \item \textbf{Systemweiten Standby und Ruhezustand komplett deaktivieren:}\\
    Da der Laptop dauerhaft als Server betrieben wird, werden sämtliche Energiespar-Targets auf Systemebene blockiert:
    \begin{itemize}
        \item \texttt{sudo systemctl mask sleep.target suspend.target hibernate.target hybrid-sleep.target}
    \end{itemize}
    Dadurch wird verhindert, dass das Betriebssystem in den Schlafmodus wechselt.
\end{enumerate}

\begin{InfoBox}{Energiespar-Sperre aufheben}
    \scriptsize Falls die Blockierung der Energiespar-Targets später wieder aufgehoben werden soll, kann dies mit dem Befehl \texttt{sudo systemctl unmask sleep.target suspend.target hibernate.target hybrid-sleep.target} rückgängig gemacht werden.
\end{InfoBox}

\section{Docker Installation}

Docker stellt eine Container-Laufzeitumgebung bereit, in der Home Assistant und weitere Dienste voneinander getrennt ausgeführt werden können. Für dieses Handbuch wird die Docker Engine direkt auf Xubuntu installiert. Docker Desktop ist für diesen Serveraufbau nicht vorgesehen, da es eine eigene Desktop-Umgebung voraussetzt.

\begin{InfoBox}{Voraussetzung für die Containerumgebung}
    \scriptsize Verwende eine Docker Engine ab Version 23.0.0. Die Installation erfolgt über die Xubuntu-Paketverwaltung und umfasst Docker Compose v2. Die benötigte Version kann nach der Installation mit \texttt{docker --version} geprüft werden.
\end{InfoBox}

\begin{enumerate}
    \item Aktualisiere zunächst die Paketquellen und installiere die Docker Engine gemeinsam mit Docker Compose v2.
    \item Aktiviere den Docker-Dienst, damit er direkt gestartet wird und künftig automatisch mit dem System hochläuft.
    \item Ergänze den aktuell angemeldeten Benutzer in der Docker-Gruppe. Dadurch können Containerbefehle später ohne \texttt{sudo} ausgeführt werden.
\end{enumerate}

\CodeFile[Installation von Docker Engine und Docker Compose]{bash}{code/docker-install.sh}\label{lst:docker-install}

\begin{InfoBox}{Änderung der Gruppenrechte übernehmen}
    \scriptsize Die neue Gruppenzugehörigkeit wird erst nach einer erneuten Anmeldung wirksam. Melde dich deshalb nach der Installation einmal ab und wieder an. Anschließend lässt sich die Installation beispielsweise mit \texttt{docker run hello-world} überprüfen.
\end{InfoBox}

\section{Router Setup und Netzstruktur}

Das Projekt-Netzwerk wird bewusst vom übrigen Netzwerk getrennt. Der Xubuntu-Server übernimmt dabei die Funktion des Gateways. Er bezieht den Internetzugang über WLAN und stellt über seine Ethernet-Schnittstelle ein eigenes lokales Netz bereit. Der verwendete A1-Router wird nicht mehr als Router oder DHCP-Server eingesetzt, sondern dient nur noch als Switch und Access Point.

\Sidebar{Aufbau des Projekt-Netzes}
Der Xubuntu-Server verbindet zwei Netze miteinander. Über WLAN erreicht er das Internet, über die Ethernet-Schnittstelle stellt er das lokale Netz \texttt{10.0.0.0/24} bereit. Die Adresse \texttt{10.0.0.1} ist dabei das Gateway für die angeschlossenen Geräte und gleichzeitig die Adresse des Servers selbst.

\subsection{A1-Router vorbereiten}

Bevor der Router mit dem Server verbunden wird, muss sein eigener Netzwerkdienst deaktiviert werden. Andernfalls könnten der Router und der Xubuntu-Server gleichzeitig IP-Adressen vergeben und es käme zu Konflikten.

\begin{enumerate}
    \item Verbinde einen Rechner mit dem A1-Router und öffne dessen Weboberfläche unter \texttt{10.0.0.138}. \Sidebar{Abweichung je nach Hersteller}
    \item Deaktiviere den DHCP-Server im Bereich für das Heimnetzwerk bzw. LAN.
    \item Entferne vorhandene statische IP-Bindungen oder Reservierungen.
    \item Belasse die Routeradresse bei \texttt{10.0.0.138}. Sie liegt damit weiterhin im später verwendeten lokalen Subnetz und bleibt erreichbar.
\end{enumerate}

\subsection{Internetfreigabe auf dem Xubuntu-Server}

Verbinde den Server zunächst über WLAN mit dem Internet. Für die Ethernet-Verbindung wird anschließend der NetworkManager-Modus \textit{shared} verwendet. In diesem Modus übernimmt der Xubuntu-Server die lokale Adressvergabe und stellt den angeschlossenen Geräten den eigenen Internetzugang über die WLAN-Verbindung bereit.

\CodeFile[Internet-Sharing und lokales Subnetz auf dem Xubuntu-Host]{bash}{code/router-network-share.sh}\label{lst:router-network-share}

Die Befehle im Listing~\ref{lst:router-network-share} führen die Konfiguration in vier Schritten durch:

\begin{enumerate}
    \item Mit \texttt{nmcli connection show} wird der Name der Ethernet-Verbindung ermittelt. Je nach Installation kann er beispielsweise \texttt{Wired connection 1} lauten.
    \item Die Verbindung erhält die Adresse \texttt{10.0.0.1/24} und wird in den gemeinsamen NetworkManager-Modus versetzt.
    \item Durch das erneute Starten der Verbindung werden die Änderungen übernommen.
    \item Mit \texttt{ip a} wird geprüft, ob die Ethernet-Schnittstelle die Adresse \texttt{10.0.0.1} verwendet.
\end{enumerate}

\begin{InfoBox}{Geräte mit dem Projekt-Netz verbinden}
    \scriptsize Verbinde den Server per LAN mit LAN-Port 1 des A1-Routers und verbinde danach die Smart-Home-Geräte per WLAN mit dem Router. Starte die Endgeräte anschließend neu, damit sie ihre Adresse vom Xubuntu-Server beziehen. Home Assistant ist nach erfolgreicher Installation (siehe Kapitel~\ref{chap:software-setup} \textit{\nameref{chap:software-setup}}) im lokalen Netz über \texttt{http://10.0.0.1:8123} erreichbar.
\end{InfoBox}

\section{Tailscale für Fernzugriff}

Das lokale Projekt-Netz ist absichtlich nicht direkt aus dem Internet erreichbar. Für die Fernwartung wird deshalb Tailscale als Overlay-Netzwerk eingesetzt. Damit kann ein autorisiertes Privatgerät eine verschlüsselte Verbindung zum Xubuntu-Server herstellen, ohne Ports im Schul- oder Heimnetz des Routers freigeben zu müssen, was aufgrund von Sicherheitsrichtlinien oft ohnehin nicht möglich ist.

\KeyImage[]{Screenshot}{Tailscale herunterladen}{../../media/Tailscale Download.png}

\begin{enumerate}
    \item Installiere Tailscale auf dem Xubuntu-Server und starte anschließend die Anmeldung.
    \item Öffne den im Terminal ausgegebenen Link und authentifiziere den Server im vorgesehenen Tailscale-Konto.
    \item Installiere Tailscale zusätzlich auf dem privaten Laptop oder Smartphone und melde auch dieses Gerät im selben Konto an.
    \item Ermittle die Tailscale-IP des Xubuntu-Servers und rufe Home Assistant (nach erfolgreicher Anmeldung) über diese Adresse und den Port \texttt{8123} auf.
\end{enumerate}

\CodeFile[Installation und Anmeldung von Tailscale auf dem Server]{bash}{code/tailscale-install.sh}

Die Anmeldung und der Verbindungsstatus können anschließend über die Tailscale-Oberfläche bzw. mit \texttt{tailscale status} kontrolliert werden. Die konkrete Tailscale-IP ist kontospezifisch und kann sich ändern, sollte daher nicht fest in irgendwelchen Konfigurationen hinterlegt werden.

\KeyImage[width=0.8\textwidth]{Screenshot}{Tailscale Pricing}{../../media/Tailscale Config 4.png}

\subsection{Technische Funktionsweise und Datenschutz}

Tailscale verwendet das WireGuard-Protokoll für die verschlüsselte Kommunikation zwischen den angemeldeten Geräten. Die zentrale Steuerungsebene vermittelt Identitäten, öffentliche Schlüssel und Verbindungsinformationen. Die Nutzdaten werden nach Möglichkeit direkt zwischen den Endgeräten übertragen. \Sidebar{Ende-zu-Ende verschlüsselt}

\begin{InfoBox}{Direkte Verbindung und Relay-Fallback}
    \scriptsize Für den Verbindungsaufbau versucht Tailscale zunächst, eine direkte Verbindung durch NAT-Traversal herzustellen. Wenn die beteiligten Netzwerke dies nicht zulassen, kann ein verschlüsselter DERP-Relay als Vermittler verwendet werden. Auch in diesem Fall bleibt der eigentliche Datenverkehr Ende-zu-Ende verschlüsselt.
\end{InfoBox}

Die privaten Schlüssel verbleiben auf den jeweiligen Endgeräten. Der Tailscale-Koordinationsdienst verwaltet daher die für den Verbindungsaufbau notwendigen Metadaten, kann den verschlüsselten Nutzdatenstrom jedoch nicht entschlüsseln. Der Zugriff auf Home Assistant ist zusätzlich auf Geräte beschränkt, die im eigenen Tailscale-Netz autorisiert wurden.

Standardmäßig wird dabei kein vollständiger Internetverkehr durch den VPN-Tunnel geleitet. Für dieses Projekt ist nur der Zugriff auf den Server und gegebenenfalls auf das Projekt-Netz erforderlich.
